\documentclass[11pt,aspectratio=169]{beamer}

\usetheme{Madrid}
\usecolortheme{beaver}
\setbeamertemplate{navigation symbols}{}

\usepackage[utf8]{inputenc}
\usepackage[T1]{fontenc}
\usepackage{amsmath,amssymb}
\usepackage{booktabs}
\usepackage{listings}
\usepackage{xcolor}
\usepackage{tcolorbox}
\tcbuselibrary{skins,breakable}

\definecolor{codebg}{RGB}{245,245,245}
\definecolor{codekw}{RGB}{0,0,180}
\lstdefinestyle{latexcode}{
  backgroundcolor=\color{codebg},
  basicstyle=\ttfamily\scriptsize,
  keywordstyle=\color{codekw}\bfseries,
  commentstyle=\color{gray},
  frame=single,
  rulecolor=\color{gray!40},
  breaklines=true,
  showstringspaces=false,
  xleftmargin=4pt,
  xrightmargin=4pt,
}
\lstset{style=latexcode}

\newtcolorbox{outputbox}{
  colback=white, colframe=black!65, boxrule=0.5pt, arc=1pt,
  left=8pt, right=8pt, top=5pt, bottom=4pt,
  title=\textsc{Compiled Output}, fonttitle=\bfseries\scriptsize,
  coltitle=white, colbacktitle=black!65
}
\newcommand{\whybox}[1]{%
  \vspace{2pt}
  {\scriptsize\raggedright\textcolor{red!55!black}{\textbf{Why \& What:}}~#1\par}
}

\newtcolorbox{exercisebox}{
  colback=orange!8, colframe=orange!70!black, boxrule=0.6pt,
  arc=2pt, left=6pt, right=6pt, top=4pt, bottom=4pt,
  title={Hands-on Exercise}, fonttitle=\bfseries, coltitle=white,
  colbacktitle=orange!70!black
}
\newtcolorbox{homeworkbox}{
  colback=blue!6, colframe=blue!50!black, boxrule=0.6pt,
  arc=2pt, left=6pt, right=6pt, top=4pt, bottom=4pt,
  title={Homework}, fonttitle=\bfseries, coltitle=white,
  colbacktitle=blue!50!black
}

\title[\texttt{\textbackslash begin} to \texttt{\textbackslash end}]{From \texttt{\textbackslash begin} to \texttt{\textbackslash end}}
\subtitle{The Complete LaTeX Guide \\[2pt] \normalsize Day 4 of 7: Citations \& Bibliography}
\author[A.\ Rahman]{Atikur Rahman}
\institute[White Globe Initiative]{White Globe Initiative}

\setbeamertemplate{footline}{
  \leavevmode
  \hbox{%
  \begin{beamercolorbox}[wd=.38\paperwidth,ht=2.5ex,dp=1.125ex,leftskip=1em]{author in head/foot}%
    \usebeamerfont{author in head/foot}White Globe Initiative
  \end{beamercolorbox}%
  \begin{beamercolorbox}[wd=.42\paperwidth,ht=2.5ex,dp=1.125ex,center]{title in head/foot}%
    \usebeamerfont{title in head/foot}From \texttt{\textbackslash begin} to \texttt{\textbackslash end} --- Day 4
  \end{beamercolorbox}%
  \begin{beamercolorbox}[wd=.20\paperwidth,ht=2.5ex,dp=1.125ex,right,rightskip=1em]{date in head/foot}%
    \usebeamerfont{date in head/foot}\insertframenumber{} / \inserttotalframenumber
  \end{beamercolorbox}}%
  \vskip0pt
}

\begin{document}

\begin{frame}[plain]
  \begin{center}
    \vspace{4pt}
    {\Huge\bfseries\usebeamercolor[fg]{title}From \texttt{\textbackslash begin} to \texttt{\textbackslash end}}\\[6pt]
    {\Large\itshape The Complete LaTeX Guide}\\[18pt]
    {\large\bfseries Day 4 of 7 --- Citations \& Bibliography}\\[22pt]
    \rule{0.7\textwidth}{0.4pt}\\[14pt]
    \begin{tabular}{rl}
      \textbf{Instructor:}       & Atikur Rahman \\[3pt]
      \textbf{Organized by:}     & White Globe Initiative \\[3pt]
      \textbf{Instructor Email:} & \texttt{atik@whiteglobeinitiative.org} \\[3pt]
      \textbf{General Contact:}  & \texttt{contact@whiteglobeinitiative.org} \\[3pt]
      \textbf{Format:}           & 7-Day Intensive Course \\[3pt]
      \textbf{Session:}          & Day 4 of 7 \\
    \end{tabular}\\[18pt]
    \rule{0.7\textwidth}{0.4pt}
  \end{center}
\end{frame}

\begin{frame}{Today's Agenda}
  \tableofcontents
\end{frame}

\section{Recap \& Today's Goal}

\begin{frame}{Quick Recap: Day 3}
  \small
  Yesterday: \texttt{tabular} + \texttt{booktabs} for
  publication-quality tables, \texttt{multicolumn}/\texttt{cmidrule}
  for grouped regression headers, figures with
  \texttt{includegraphics}, float placement, and the
  tablesgenerator.com workflow.
  \vspace{6pt}

  \textbf{Today's goal:} economics journals almost universally use
  author-year citation style. Leave today with a working BibTeX
  pipeline you can reuse for every future paper -- we'll practice on
  real published research from White Globe Initiative's own team.
\end{frame}

\section{BibTeX Basics}

\begin{frame}[fragile]{The .bib File}
  \begin{lstlisting}
@article{nandi2022financial,
  author  = {Nandi, Biplob Kumar and Hasan, Gazi Quamrul
             and Kabir, Mohammad Habibullah},
  title   = {A Tale of the Financial Inclusion-Growth Nexus},
  journal = {Journal of Financial Economic Policy},
  year    = {2022},
  volume  = {14},
  number  = {3},
  pages   = {381--402}
}
  \end{lstlisting}
  \whybox{A separate \texttt{references.bib} file holds all your
  sources -- one \texttt{@}-entry per source, with \textbf{no visible
  output by itself}: it's raw data, rendered only when you cite it
  (next slides). Entry types: \texttt{@article}, \texttt{@book},
  \texttt{@incollection}, \texttt{@techreport}. The key
  (\texttt{nandi2022financial}) is what you'll type to cite this
  source in your text -- pick something memorable: author + year + one
  keyword.}
\end{frame}

\section{natbib \& Author-Year Styles}

\begin{frame}[fragile]{natbib: Author-Year Citations}
  \begin{lstlisting}
\usepackage[authoryear]{natbib}
\bibliographystyle{agsm}   % or apalike -- common in economics
...
\bibliography{references}
  \end{lstlisting}
  \begin{outputbox}
    \footnotesize
    Nandi, B.K., Hasan, G.Q., and Kabir, M.H. (2022). A tale of the
    financial inclusion-growth nexus. \textit{Journal of Financial
    Economic Policy}, 14(3), 381--402.
  \end{outputbox}
  \whybox{\texttt{natbib} is the standard package for author-year
  citation styles used by most economics journals. \texttt{agsm} and
  \texttt{apalike} are common bibliography styles -- journal templates
  usually specify their own (Day 6). \texttt{\textbackslash
  bibliography\{references\}} points to your \texttt{.bib} file (no
  extension) and is what actually \textbf{renders} the reference-list
  entry shown above, formatted automatically from the raw data you
  typed on the previous slide.}
\end{frame}

\begin{frame}[fragile]{Citation Commands: citet and citep}
  \begin{lstlisting}
\citet{nandi2022financial} find a positive link between
financial inclusion and growth.

Financial deepening supports growth
\citep{nandi2022financial}.
  \end{lstlisting}
  \begin{outputbox}
    \small
    Nandi et al.\ (2022) find a positive link between financial
    inclusion and growth.\\[6pt]
    Financial deepening supports growth (Nandi et al., 2022).
  \end{outputbox}
  \whybox{\texttt{\textbackslash citet\{...\}} -- ``\textbf{t}extual'':
  the authors' names become part of the sentence, exactly like ``Nandi
  et al.\ (2022) find...'' above. \texttt{\textbackslash
  citep\{...\}} -- ``\textbf{p}arenthetical'': the citation sits in
  parentheses instead. Rule of thumb: if the authors are the subject of
  your sentence, use \texttt{citet}; otherwise \texttt{citep}. Note
  \texttt{natbib} auto-inserts ``et al.'' for three or more authors.}
\end{frame}

\section{biblatex: A Brief Note}

\begin{frame}{biblatex as an Alternative}
  \small
  \texttt{biblatex} (with the \texttt{biber} backend) is a more modern,
  flexible bibliography system that supports more citation styles out
  of the box and handles non-English references better. For this
  course, \textbf{natbib is sufficient} -- it is what most economics
  journal templates already expect, so it is the lower-friction choice
  while you are building the habit. You can switch to
  \texttt{biblatex} later without relearning your \texttt{.bib} files
  -- the entries themselves don't change, only the packages that read
  them.
\end{frame}

\section{Sourcing .bib Entries}

\begin{frame}{Where .bib Entries Come From}
  \small
  \begin{itemize}
    \item \textbf{Google Scholar}: search the paper $\rightarrow$ click the quotation-mark ``Cite'' icon $\rightarrow$ \texttt{BibTeX} link $\rightarrow$ copy-paste.
    \item \textbf{Reference managers} (Zotero, Mendeley): export your library directly as \texttt{.bib}.
    \item \textbf{Journal websites}: many have a direct ``Export citation $\rightarrow$ BibTeX'' button.
    \item \textbf{Practice source for today}: browse
      \texttt{whiteglobeinitiative.org/published-papers} -- build real
      \texttt{.bib} entries from WGI's own peer-reviewed output.
  \end{itemize}
  \vspace{4pt}
  Always spot-check auto-generated entries -- author name formatting
  and page ranges are the most common errors.
\end{frame}

\section{Cross-Referencing}

\begin{frame}[fragile]{Cross-Referencing Sections, Equations, Tables, Figures}
  \begin{lstlisting}
\section{Empirical Strategy}
\label{sec:strategy}

As discussed in Section~\ref{sec:strategy}, we estimate...
  \end{lstlisting}
  \begin{outputbox}
    \small As discussed in Section~2, we estimate...
  \end{outputbox}
  \whybox{Same \texttt{\textbackslash label} / \texttt{\textbackslash
  ref} mechanism from Day 2 (equations) and Day 3 (tables/figures) --
  now applied to sections. \texttt{\textbackslash ref} outputs just the
  number; for equations, \texttt{\textbackslash eqref} additionally
  wraps it in parentheses. These numbers update automatically as you
  add or reorder content -- never hand-count section numbers again.}
\end{frame}

\section{Practice}

\begin{frame}{Hands-on Exercise}
  \begin{exercisebox}
    \begin{enumerate}
      \item Visit \texttt{whiteglobeinitiative.org/published-papers}
        and build a \textbf{five-entry} \texttt{.bib} file from real
        papers listed there (or your own field's literature).
      \item Write a short mock literature-review paragraph (4--5
        sentences) citing all five.
      \item Use \texttt{\textbackslash citet} at least twice and
        \texttt{\textbackslash citep} at least twice.
      \item Compile and confirm your reference list appears correctly
        at the end of the document.
    \end{enumerate}
  \end{exercisebox}
\end{frame}

\begin{frame}{Homework Before Day 5}
  \begin{homeworkbox}
    Build out the real bibliography for your running paper draft and
    insert citations throughout the introduction / literature review
    section. Aim for at least 8--10 sources.
  \end{homeworkbox}
  \vspace{10pt}
  \small\textbf{Resources for today:}
  \begin{itemize}\small
    \item Overleaf -- ``Bibliography management with bibtex'' / ``with natbib''
    \item Overleaf -- ``Natbib citation styles''
    \item Overleaf -- ``Cross referencing sections, equations and floats''
    \item Wikibooks LaTeX -- Bibliography Management chapter
  \end{itemize}
\end{frame}

\begin{frame}{Coming Up: Day 5}
  \Large
  \centering
  Document Structure \& Long-Document Tools \\[6pt]
  \normalsize
  Table of contents, footnotes, hyperlinks, and splitting a long paper
  across multiple files.
\end{frame}

\begin{frame}[plain]
  \begin{center}
    \vspace{35pt}
    {\Large\bfseries Thank You}\\[10pt]
    {\large From \texttt{\textbackslash begin} to \texttt{\textbackslash end}: The Complete LaTeX Guide}\\[4pt]
    Day 4 of 7 --- Citations \& Bibliography\\[20pt]
    \rule{0.5\textwidth}{0.4pt}\\[10pt]
    Atikur Rahman \\
    White Globe Initiative \\[2pt]
    \texttt{atik@whiteglobeinitiative.org} \\
    \texttt{contact@whiteglobeinitiative.org}
  \end{center}
\end{frame}

\end{document}
